\documentclass{subfiles}
\begin{document}
    \begin{figure}[!hb]
        \centering
        \begin{subfigure}{0.9\textwidth}
            \centering
            \begin{tikzpicture}
                [
                    internal/.style = { draw, circle, minimum size = 0.1cm, fill=rpText },
                    subtree/.style = { draw, regular polygon, regular polygon sides = 3, minimum size = 0.25cm, fill = rpText },
                    every label/.style = {font = \ttfamily\tiny}
                ]

                \graph [tree layout, math nodes] {
                    "" [internal, "$x.p$" left] -- {
                        "" [internal, "$x$" left] -- {
                            "" [subtree, "$\alpha$" left],
                            "" [subtree, "$\beta$" right]
                        },
                        "" [subtree, "$\gamma$" right]
                    }
                };

                \draw [-Stealth, thick] (1.5, -1) -- (2.5, -1);

                \graph [tree layout, math nodes] {
                    "" [internal, "$x$" left, desired at={(4, 0)}] -- {
                        "" [subtree, "$\alpha$" left],
                        "" [internal, "old $x.p$" right] -- {
                            "" [subtree, "$\beta$" left],
                            "" [subtree, "$\gamma$" right]
                        }
                    }
                };

            \end{tikzpicture}
            \caption{}\label{subfig:splaying-case-1}
        \end{subfigure}
        \vskip 2.5pt
        \begin{subfigure}{0.9\textwidth}
            \centering
            \begin{tikzpicture}
                [
                    internal/.style = { draw, circle, minimum size = 0.1cm, fill=rpText },
                    subtree/.style = { draw, regular polygon, regular polygon sides = 3, minimum size = 0.25cm, fill = rpText },
                    every label/.style = {font = \ttfamily\tiny}
                ],

                \graph [tree layout] {
                    "" [internal, "$x.p.p$" left] -- {
                        "" [internal, "$x.p$" left] -- {
                            "" [internal, "$x$" left] -- { 
                                "" [subtree, "$\alpha$" left],
                                "" [subtree, "$\beta$" right]
                            },
                            "" [subtree, "$\gamma$" right]
                        },
                        "" [subtree, "$\delta$" right]
                    }
                };

                \draw [-Stealth, thick] (1, -2) -- (2, -2);

                \graph [tree layout] {
                    "" [internal, "$x.p$" left, desired at={(3.5, -1)}] -- {
                        "" [internal, "$x$" left] -- { 
                            "" [subtree, "$\alpha$" left],
                            "" [subtree, "$\beta$" left]
                        },
                        "" [internal, "old $x.p.p$" left] -- {
                            "" [subtree, "$\gamma$" left],
                            "" [subtree, "$\delta$" left]
                        }
                    }
                };

                \draw [-Stealth, thick] (5, -2) -- (6, -2);

                \graph [tree layout] {
                    "" [internal, "$x$", desired at={(7, 0)}] -- {
                        "" [subtree, "$\alpha$" left],
                        "" [internal, "old $x.p$" right] -- {
                            "" [subtree, "$\beta$" left],
                            "" [internal, "old $x.p.p$" right] -- {
                                "" [subtree, "$\gamma$" left],
                                "" [subtree, "$\delta$" right]
                            }
                        }
                    }
                };
            \end{tikzpicture}
            \caption{}\label{subfig:splaying-case-2}
        \end{subfigure}
        \vskip 2.5pt 
        \begin{subfigure}{0.9\textwidth}
            \begin{tikzpicture}
                                [
                    internal/.style = { draw, circle, minimum size = 0.1cm, fill=rpText },
                    subtree/.style = { draw, regular polygon, regular polygon sides = 3, minimum size = 0.25cm, fill = rpText },
                    every label/.style = {font = \ttfamily\tiny}
                ],

                \graph [tree layout] {
                    "" [internal, "$x.p.p$" left] -- {
                        "" [internal, "$x.p$" left] -- {
                            "" [subtree, "$\gamma$" left],
                            "" [internal, "$x$" left] -- { 
                                "" [subtree, "$\alpha$" left],
                                "" [subtree, "$\beta$" right]
                            }
                        },
                        "" [subtree, "$\delta$" right]
                    }
                };

                \draw [-Stealth, thick] (1, -1.5) -- (2, -1.5);

                \graph [tree layout] {
                    "" [internal, "old $x.p.p$" left, desired at={(4, 0)}] -- {
                        "" [internal, "$x$" left] -- {
                            "" [internal, "old $x.p$" left] -- {
                                "" [subtree, "$\gamma$" left],
                                "" [subtree, "$\alpha$" right]
                            },
                            "" [subtree, "$\beta$" right]
                        },
                        "" [subtree, "$\delta$" right]
                    }
                };

                \draw [-Stealth, thick] (5, -1.5) -- (6, -1.5);

                \graph [tree layout] {
                    "" [internal, "$x$" left, desired at={(7.5, -1)}] -- {
                        "" [internal, "old $x.p$" left] -- {
                            "" [subtree, "$\gamma$" left],
                            "" [subtree, "$\alpha$" left]
                        },
                        "" [internal, "old $x.p.p$" right] -- {
                            "" [subtree, "$\beta$" left],
                            "" [subtree, "$\delta$" left]
                        }
                    }
                };
            \end{tikzpicture}
            \caption{}\label{subfig:splaying-case-3}
        \end{subfigure}
        \caption{Example of splaying steps. Circles denotes internal nodes, 
            whilst triangles denote arbitrary subtrees.
            \ref{subfig:splaying-case-1} shows the case in which $x.p$ is the root.
            As stated, a simple rotation to the right will suffice.
            \ref{subfig:splaying-case-2} shows the case in which both $x$ and $x.p$ are left children;
            a left rotation on $x.p$ first, and a right one on $x$ suffice.
            \ref{subfig:splaying-case-3} shows the case in which $x$ is a right child, 
            and $x.p$ a left child; we rotate on $x$ twice.
        }\label{fig:splay-example-splaying}
    \end{figure}
\end{document}
